\documentclass[11pt,twoside,a4paper]{article}

\usepackage{rotating,booktabs,graphicx}

\title{HSCP for Phase2 Trk TDR}

\begin{document}

\section{Introduction}

Several theoretical extensions of the Standard Model (SM), including Super Symmetry (split SUSY and MSSM), Gauge mediated Symmetry Breaking (GMSB), Hidden-Valley, Extra dimensions, and several others, predict the existence of heavy particles with long lifetime.
If produced with lifetime larger that a few nanoseconds, these particles can pass through part of the detector before decaying, and therefore appear as stable.
Depending on the mass and the charge of the HSCP particle, anomalously high rates of energy losses through ionization ($\mathrm{d}E/\mathrm{d}x$) with respect to the typical energy loss for SM particles ($\approx 3 \mathrm{MeV}/\mathrm{cm}$ for SM particles with 10-1000 GeV momentum) can be expected. 
This feature is exploited to search for heavy stable (or quasi-stable) charged particles (HSCPs) by measuring the energy loss in the CMS silicon Tracker, discriminating heavily ionizing particles (HIPs) from minimum ionizing particles (MIPs).

\section{Samples}

Two benchmark scenarios are employed to account for a wider range of possible signatures of HSCP production at LHC.

A first kind of signal considered in this study is the production of gluino ($\tilde{g}$) pairs, that interact via the strong, electromagnetic, and weak interactions. 
In this model, HSCPs can also bound with SM quarks to form R-hadrons by strong interaction, leading to modifications of their quark constituents and therefore of their electric charge during the propagation through the detector material.
Signal gluino have been generated with masses in the range 600-2600 GeV under the Split SUSY scenarios under the high squark mass (10 TeV) assumption, by the PYTHIA v8.xxx generator with the tune CUETP8M1.
The fraction $f$ of produced $\tilde{g}$ hadronizing into a $\tilde{g}$-g state (R-gluonball) is an unknown parameter of the hadronization model and affects the fraction of R-hadrons that are neutral at production, and for this study results are obtained under the $f = 0.1$ assumption.

The second kind of signal consists of lepton-like HSCPs, which interacts via electromagnetic and weak forces behaving similarly to muons, as they tend to be deeply penetrating into the detector.
The benchmark chosen for lepton-like HSCP is the minimal gauge mediated supersymmetry breaking (mGMSB) model, where production of supersymmetric quasi-stable leptons ($\tilde{\tau}$) can proceed either directly or via production of heavier supersymmetric particles (mainly squarks and gluino pairs) further decaying to one or more $\tilde{\tau}$ particles.
The mGMSB model is chosen as benchmark using the SPS7 slope, which has the stau as the next-to-lightest supersymmetric particle (NLSP). 
The particle mass spectrum and the decay table are produced with the program ISASUGRA version 7.69. The mGMSB parameter $\Lambda$ is varied from 31 to
160 TeV, with fixed parameters $\mathrm{N}_{\mathrm{mes}} = 3$, $\tan \beta = 10$, $\mu > 0$, $\mathrm{C}_{\mathrm{grav}} = 10000$, and $\mathrm{M}_{\mathrm{mes}}/\Lambda = 2$. 
The large value of $\mathrm{C}_{\mathrm{grav}}$ results in a long-lived stau, while $\Lambda = 31-510$ TeV gives a stau mass of 200 to 1600 GeV. 
Signal is generated using PYTHIA v8.xxx with the default tune CUETP8M1.

SM background samples have been used to simulate MIP tracks produced as the result of proton-proton interactions in the CMS detector.
For this purpose, Minimum Bias samples have been used, generated assuming a proton-proton inelastic cross section of xxx mb.

To simulate high transverse momentum muons traversing the detector, a samples of Drell-Yan production is also used, to exploit muons from the leptonic Z boson decay, as well as leptonic decaying top quark pairs.
Drell-Yan samples have been produced with the MadgraphMLM generator interfaced with PYTHIA to describe the hadronization, whereas the $\mathrm{t}\bar{\mathrm{t}}$ SM background the Powheg generator has been used, also in this case interfaced with PYTHIA8.

All the signal and background samples used in this analysis have been subjected to the full simulation of the CMS detector with the phase 2 geometry (2023-D4) using GEANT 4. 
For all signal and background samples, three different configuration have been produced to account for additional events produced in association with the primary collision in the same beam crossing (pileup): no pileup, 140 pileup, and 200 pileup.

\begin{sidewaystable}
\small
\centering
\caption{Simulated samples used for this study. All the samples have been produced in three different pileup configurations: NoPU, 140PU, and 200PU.\label{tab:samples}}
\begin{tabular}{ll}
\toprule
Sample & $\sigma$ \\
\midrule
\verb|/HSCPgluino_M_600_TuneCUETP8M1_14TeV_pythia8/PhaseIIFall16DR82-*_90X_upgrade2023_realistic_v1-v1/GEN-SIM-RECO|	& xxx \\
\verb|/HSCPgluino_M_1000_TuneCUETP8M1_14TeV_pythia8/PhaseIIFall16DR82-*_90X_upgrade2023_realistic_v1-v1/GEN-SIM-RECO|	& xxx \\
\verb|/HSCPgluino_M_1800_TuneCUETP8M1_14TeV_pythia8/PhaseIIFall16DR82-*_90X_upgrade2023_realistic_v1-v1/GEN-SIM-RECO|	& xxx \\
\verb|/HSCPgluino_M_1400_TuneCUETP8M1_14TeV_pythia8/PhaseIIFall16DR82-*_90X_upgrade2023_realistic_v1-v1/GEN-SIM-RECO|	& xxx \\
\verb|/HSCPgluino_M_2600_TuneCUETP8M1_14TeV_pythia8/PhaseIIFall16DR82-*_90X_upgrade2023_realistic_v1-v2/GEN-SIM-RECO|	& xxx \\
\verb|/HSCPgluino_M_2200_TuneCUETP8M1_14TeV_pythia8/PhaseIIFall16DR82-*_90X_upgrade2023_realistic_v1-v1/GEN-SIM-RECO|	& xxx \\
\midrule
\verb|/HSCPppstau_M_200_TuneCUETP8M1_14TeV_pythia8/PhaseIIFall16DR82-*_90X_upgrade2023_realistic_v1-v1/GEN-SIM-RECO|	& xxx \\
\verb|/HSCPppstau_M_432_TuneCUETP8M1_14TeV_pythia8/PhaseIIFall16DR82-*_90X_upgrade2023_realistic_v1-v1/GEN-SIM-RECO|	& xxx \\
\verb|/HSCPppstau_M_651_TuneCUETP8M1_14TeV_pythia8/PhaseIIFall16DR82-*_90X_upgrade2023_realistic_v1-v1/GEN-SIM-RECO|	& xxx \\
\verb|/HSCPppstau_M_871_TuneCUETP8M1_14TeV_pythia8/PhaseIIFall16DR82-*_90X_upgrade2023_realistic_v1-v1/GEN-SIM-RECO| 	& xxx \\
\verb|/HSCPppstau_M_1218_TuneCUETP8M1_14TeV_pythia8/PhaseIIFall16DR82-*_90X_upgrade2023_realistic_v1-v1/GEN-SIM-RECO|	& xxx \\
\verb|/HSCPppstau_M_1599_TuneCUETP8M1_14TeV_pythia8/PhaseIIFall16DR82-*_90X_upgrade2023_realistic_v1-v1/GEN-SIM-RECO|	& xxx \\
\midrule
\verb|/MinBias_*_TuneCUETP8M1_14TeV-pythia8/PhaseIIFall16DR82-*_90X_upgrade2023_realistic_v1-v1/GEN-SIM-RECO|		& xxx \\
\midrule
\verb|/DYJetsToLL_M-50_TuneCUETP8M1_14TeV-madgraphMLM-pythia8_ext1/PhaseIIFall16DR82-*_90X_upgrade2023_realistic_v1-v1/GEN-SIM-RECO|	& xxx \\
\verb|/TTTo2L2Nu_TuneCUETP8M1_14TeV-powheg-pythia8/PhaseIIFall16DR82-*_90X_upgrade2023_realistic__v1_ext1-v1/GEN-SIM-RECO|	& xxx \\
\bottomrule
\end{tabular}
\end{sidewaystable}

\section{Event Selection}

Events are selected by requiring the presence of at least one muon with high transverse momentum or an high missing transverse momentum $E_{\mathrm{T}}^{\mathrm{miss}}$ as reconstructed from the particle-flow algorithm, to emulate the typical trigger requirements applied for this analysis.
Single muon triggers are especially indicated for lepton-like HSCPs, which interact weakly with the detector and thus are normally reconstructed as muons inside the detector.
To recover some efficiency for hadron-like HSCPs, for which R-hadrons can change charge as they propagate through the detector and become neutral within or outside the tracker, online selection criteria relying on the missing transverse momentum reconstructed in the detector are also exploited to complement the single muon triggers.
For the purpose of this study, online trigger performances similar to the ones of Run2 have been assumed, and events are retained only if fulfilling the offline selection cuts of having at least one muon with $p_{\mathrm{T}} > 55$ GeV or $E_{\mathrm{T}}^{\mathrm{miss}} > 200$ GeV, values for which the trigger is assumed to be fully efficient. 

Candidates HSCP are considered all the tracks with $p_{\mathrm{T}} > 55$ GeV (as measured in the tracker detector only) in the tracker detector accept, with a relative uncertainty on $p_{\mathrm{T}}$ less than 25\%, pseudorapidity in the tracker acceptance $|\eta| < 4$, track fit $\chi^2/\mathrm{dof} < 5$, and magnitudes of the longitudinal and transverse impact parameters with respect to the primary vertex less than 0.5 cm. 
Candidates tracks must have at least XXX measurements in the silicon pixel detector.

\begin{table}
 \centering
 \caption{Selection criteria\label{tab:selection}}
 \begin{tabular}{ll}
  \toprule
  Variable 			& cut \\ 
  \midrule	
  $p_{\mathrm{T}}$		& $> 55$ GeV \\
  $|\eta|$			& $< 4$ \\
  $d_z$ and $d_{xy}$		& $< 0.5$ cm \\
  $\sigma(p_{\mathrm{T}})/p_{\mathrm{T}}$		& $< 0.25$ \\
  track $\chi^2/$dof		& $> 5$ \\
  pixel hits			& $> X$ \\
  strip hits			& $> X$ \\
  $\mathrm{d}E/\mathrm{d}x$ measurements			& $> X$ \\
  \bottomrule
 \end{tabular}
\end{table}

\begin{figure}
\centering
\includegraphics[width=.5\textwidth]{figures/dummy.pdf}
\caption{Signal efficiency versus the track pseudorapidity $\eta$.\label{fig:effeta}}
\end{figure}


\section{Ionizing Particle Discriminator}

An estimator of the degree of compatibility of the track with the hypothesis of it being generated by a MIP is defined in the analysis to separate candidate HSCP from tracks from SM background sources.
The $\mathrm{d}E/\mathrm{d}x$ discriminator $I_{\mathrm{as}}$ is define as:

$$
I_{\mathrm{as}} = \frac{3}{N} \times \left( \frac{1}{12N} + \sum_{i=1}^{N} \left[ P_i \times \left( P_i - \frac{2i-1}{2N} \right)^2 \right] \right)
$$

where $N$ is the number of measurements in the silicon-tracker detectors, $P_i$ is the probability for a minimum-ionizing particle to produce a charge smaller or equal to that of the $i$-th measurement for the observed path length in the detector, and the sum is over the track measurements ordered in terms of increasing $P_i$.

Pixel measurements are used for the computation of the discriminator due to the high resolution of the $\mathrm{d}E/\mathrm{d}x$ measurements.
Strip modules provide a coarser $\mathrm{d}E/\mathrm{d}x$ information with respect to the pixel modules: a Hit-Over-Threshold information (HoT bit) is provided only in the case an energy deposit greater that $1.4$ times the one expected from a MIP is measured.
For this reason, three different implementations of the $I_{\mathrm{as}}$ discriminator are implemented in order to investigate which strategy could exploit the most out of the PS modules information.

\begin{itemize}
 \item[i.]   Evaluate a simplified pixel-only discriminator to exploit the higher granularity of the pixel detector, and recompute $I_{\mathrm{as}}$ including the HoT bits only if the the pixel-only discriminator is above the threshold of XXX.
 \item[ii.]  Define multiple discriminator templates based on MinBias depending on the multiplicity of HoT bits associated to the track
 \item[iii.] Build different discriminators based on the  geometry type of APV modules
\end{itemize}

Receiver operating characteristic (ROC) curves have been evaluated to evaluate the strategy providing the best performances.
% Signal samples have been exploited to measure the signal efficieny, while the SM backgound sample have been used to estimate the probability o

\begin{figure}
\centering
\includegraphics[width=.5\textwidth]{figures/dummy.pdf}
\caption{$\mathrm{d}E/\mathrm{d}x$ profile versus the momentum of the track.\label{fig:dedxvsp}}
\end{figure}

\section{Background Estimation}

The HSCP candidates are selected using the $I_{\mathrm{as}}$ discriminator by requiring tracks with  criteria are $p_{\mathrm{T}} > 65$ GeV and $I_{\mathrm{as}} > 0.3$. 
The SM background contribution is predicted relying on two non-correlated variables discriminating between signal and background.
The background yield is in fact estimated from the Drell-Yan and $\mathrm{t}\bar{\mathrm{t}}$ MC samples with a "ABCD" procedure.
The number of background events estimated in the signal region (region D) is obtained by correcting the number of events passing the $I_{\mathrm{as}}$ requirement only (region B), for the ratio of the events yields passing the $p_{\mathrm{T}}$ requirement alone (region C) divided by the number of events failing both requirements (region A).

\begin{figure}
\centering
\includegraphics[width=.5\textwidth]{figures/dummy.pdf}
\caption{ROC\label{fig:roc}}
\end{figure}

\begin{table}
 \centering
 \caption{Selection criteria\label{tab:numbers}}
 \begin{tabular}{ll}
  \toprule
  region 	& n. events \\
  \midrule
  A		& X \\
  B		& X \\
  C		& X \\
  D		& X \\
  \bottomrule
 \end{tabular}
\end{table}

\section{Results}

The expected performances of the analysis in absence of signal are shown in terms of cross-section reach, defined as the cross-section for which we expect to observe signal with a significance of at least 5 standard deviations (5$\sigma$).
The results are presented for the three pileup scenarios (NoPU, 140PU, 200PU) for an integrated luminosity of 3000 $\mathrm{fb}^{-1}$.

\begin{figure}
\centering
\includegraphics[width=.5\textwidth]{figures/dummy.pdf}
\caption{Limit\label{fig:limit}}
\end{figure}


\end{document}

